\documentclass[a4paper,11pt]{article}
\usepackage[utf8]{inputenc}
\usepackage{graphicx}
\usepackage{float}
\usepackage{amsmath}
\usepackage{geometry}
\usepackage{booktabs}
\usepackage{siunitx}
\usepackage{hyperref}

\geometry{margin=1in}

\title{Laboratory Report: The Damped Driven Harmonic Oscillator \\ \large Section: Antenna Measurement}
\author{Eni Hoxhalli \& Partner Name}
\date{November 27, 2025}

\begin{document}

\maketitle

\section{Introduction}

For the last part of our lab session with the Rohde \& Schwarz FPC1500 Spectrum Analyzer, we characterized a Molex 206760 Cellular Antenna. The goal was to measure the antenna's return loss ($S_{11}$) and compare it to the spec sheet, then use it to pick up ambient signals in the lab.

Return loss ($S_{11}$) tells us how well the antenna is matched to the transmission line. It's the ratio of reflected power to incident power in dB. A more negative value means better matching and more efficient radiation. You want to see a dip at the operating frequency where power is radiated instead of reflected back.

\section{Experimental Setup}

We used a test board with a surface-mount Molex cellular antenna connected through SMA connectors. We ran two different measurement configurations:

\subsection{Return Loss Measurement (S11)}
\begin{itemize}
    \item \textbf{Configuration:} Tracking Generator (TG) mode enabled
    \item \textbf{Generator Output Power:} $-20\,\text{dBm}$
    \item \textbf{Frequency Range:} $5\,\text{kHz}$ to $3\,\text{GHz}$
    \item \textbf{Resolution Bandwidth (RBW):} Auto
    \item \textbf{Video Bandwidth (VBW):} Auto
    \item \textbf{Measurement:} Power reflected back to port 1 (S11)
\end{itemize}

\subsection{Ambient Signal Detection}
\begin{itemize}
    \item \textbf{Configuration:} Internal source turned off
    \item \textbf{Frequency Range:} $5\,\text{kHz}$ to $3\,\text{GHz}$
    \item \textbf{Trace Mode:} Maximum hold
    \item \textbf{Measurement:} External signals received by antenna
\end{itemize}

\section{Results and Analysis}

\subsection{Antenna Return Loss ($S_{11}$)}

The return loss tells us how well the antenna radiates. A dip in the measured power means resonance—energy is being radiated instead of reflected. The Molex 206760 datasheet says it's designed for cellular bands:

\begin{itemize}
    \item \textbf{Operating Range:} $1710\,\text{MHz}$ to $2700\,\text{MHz}$
    \item \textbf{Resonant Frequency:} $1834\,\text{MHz}$
    \item \textbf{Minimum Return Loss:} Around $-42.5\,\text{dB}$ at resonance
\end{itemize}

Figure~\ref{fig:s11_full} shows what we measured across the full $3\,\text{GHz}$ range.

\begin{figure}[H]
    \centering
    \includegraphics[width=0.95\textwidth]{s11_full_spectrum.png}
    \caption{Measured $S_{11}$ from $5\,\text{kHz}$ to $3\,\text{GHz}$. Baseline around $-20\,\text{dBm}$ (almost everything reflected). Green area is the spec operating range, red line is where resonance should be at $1834\,\text{MHz}$.}
    \label{fig:s11_full}
\end{figure}

\subsubsection{Comparison with Specification Sheet}

Table~\ref{tab:spec_vs_meas} compares what the manufacturer claims versus what we actually measured.

\begin{table}[H]
    \centering
    \caption{Comparison of Specification vs. Measured Data}
    \label{tab:spec_vs_meas}
    \begin{tabular}{@{}lcc@{}}
        \toprule
        \textbf{Parameter} & \textbf{Specification} & \textbf{Measured} \\
        \midrule
        Resonant Frequency & $1834\,\text{MHz}$ & $1834\,\text{MHz}$ (no dip observed) \\
        Return Loss at Resonance & $-42.5\,\text{dB}$ & $-21.3\,\text{dBm}$ \\
        Discrepancy & --- & $+21.2\,\text{dB}$ \\
        \bottomrule
    \end{tabular}
\end{table}

So here's the problem: at $1834\,\text{MHz}$, we measured $-21.3\,\text{dBm}$, which is barely different from the generator baseline of $-20\,\text{dBm}$. This means almost everything is getting reflected back—the antenna isn't radiating at all. We should've seen that deep dip down to $-42.5\,\text{dB}$, but it just wasn't there.

Figure~\ref{fig:s11_cellular} gives a closer look at the cellular band region ($1.7$--$2.7\,\text{GHz}$) where this antenna is supposed to work best.

\begin{figure}[H]
    \centering
    \includegraphics[width=0.95\textwidth]{s11_cellular_band.png}
    \caption{Cellular band ($1710$--$2700\,\text{MHz}$) zoomed in. Red line is where we should see $-42.5\,\text{dB}$ at $1834\,\text{MHz}$. No dip visible.}
    \label{fig:s11_cellular}
\end{figure}

\subsubsection{Performance Across Frequency Range}

The antenna didn't work well across most of the spectrum. From $0$--$2300\,\text{MHz}$ (including the whole cellular band it's designed for), the trace stayed flat near baseline - high impedance mismatch, basically no absorption. Above $2300\,\text{MHz}$ we saw a slight dip to about $-23.8\,\text{dBm}$ around $2380\,\text{MHz}$, but that's still way below spec.

\subsubsection{Why the Results Don't Match the Spec}

The missing resonance and poor return loss are probably due to our experimental setup:

\paragraph{Bad Connections}
The lab manual mentions the test board has "crude" SMA connections. In RF, connection quality really matters. Bad soldering or loose connectors create impedance mismatches that reflect signals before they even reach the antenna. Since we're seeing $S_{11} \approx -20\,\text{dBm}$, most of the reflection is probably happening right at the connector.

\paragraph{Environmental Effects}
The manufacturer measured this in free space, but we had our board on a metal lab bench. The bench acts like a ground plane and changes the antenna's electrical length. Other stuff around the lab (equipment, walls, materials) can shift the resonant frequency by hundreds of MHz or dampen it completely.

\paragraph{Test Setup Issues}
Our test board doesn't match how this antenna would be mounted in actual devices. It's designed for a specific PCB layout with impedance matching networks that we don't have.

\subsection{Ambient Frequency Detection}

For the second part, we turned off the internal source and used the antenna as a receiver to see what signals are in the lab. The noise floor was around $-76\,\text{dBm}$.

Figure~\ref{fig:ambient_full} shows the full ambient spectrum.

\begin{figure}[H]
    \centering
    \includegraphics[width=0.95\textwidth]{ambient_full_spectrum.png}
    \caption{Ambient spectrum from $5\,\text{kHz}$ to $3\,\text{GHz}$. Noise floor around $-76\,\text{dBm}$. The peaks are wireless signals.}
    \label{fig:ambient_full}
\end{figure}

\subsubsection{Detected Signals}

Even with the bad $S_{11}$, the antenna still picked up signals. Table~\ref{tab:detected_signals} shows what we found.

\begin{table}[H]
    \centering
    \caption{Detected Ambient Signals}
    \label{tab:detected_signals}
    \begin{tabular}{@{}cccc@{}}
        \toprule
        \textbf{Frequency} & \textbf{Power} & \textbf{SNR} & \textbf{Likely Source} \\
        \midrule
        $2403\,\text{MHz}$ & $-62.5\,\text{dBm}$ & $13.5\,\text{dB}$ & Wi-Fi / Bluetooth \\
        \bottomrule
    \end{tabular}
\end{table}

The strongest signal was at $2403\,\text{MHz}$ with $-62.5\,\text{dBm}$ received power, about $13.5\,\text{dB}$ above the noise floor.

\subsubsection{Signal Identification}

The $2403\,\text{MHz}$ signal is in the $2.4\,\text{GHz}$ ISM band ($2400$--$2485\,\text{MHz}$), which is unlicensed spectrum for Wi-Fi, Bluetooth, and other stuff.

Wi-Fi channel 1 is centered at $2412\,\text{MHz}$, so our signal at $2403\,\text{MHz}$ could be the lower edge of that. But it also matches Bluetooth channel 1 exactly—Bluetooth uses $1\,\text{MHz}$ spacing from $2402$--$2480\,\text{MHz}$. Given all the Bluetooth devices in the lab (keyboards, mice, headphones, phones), that's probably what we're seeing. Or maybe both at once, since the spectrum analyzer can't really tell them apart.

Figure~\ref{fig:ambient_ism} shows the ISM band in more detail.

\begin{figure}[H]
    \centering
    \includegraphics[width=0.95\textwidth]{ambient_ism_band.png}
    \caption{$2.4\,\text{GHz}$ ISM band ($2.3$--$2.5\,\text{GHz}$). Red circle is our peak at $2403\,\text{MHz}$. Green lines show Wi-Fi channels 1, 6, and 11.}
    \label{fig:ambient_ism}
\end{figure}

\section{Conclusion}

The antenna is supposed to work well from $1710$--$2700\,\text{MHz}$ with resonance at $1834\,\text{MHz}$, but our measurements didn't match the spec sheet at all. $S_{11}$ stayed around $-20\,\text{dBm}$ across the cellular band - basically everything got reflected. This was probably because of the crude SMA connections on the test board, plus having it sitting on a metal bench instead of free space like the manufacturer tested it. Our setup also doesn't have the matching networks this antenna needs.

Even with the bad matching, we still picked up a signal at $2403\,\text{MHz}$ that's likely Bluetooth or Wi-Fi from devices in the lab. So even a poorly matched antenna can receive if the signal is strong enough.

The main takeaway is that spec sheet performance assumes ideal conditions. In a real lab setup with imperfect connections and environmental interference, you won't get those numbers without proper equipment like a VNA with calibration.

\section*{Acknowledgments}
Thanks to the lab instructors for the help with the spectrum analyzer and for answering our questions.

\end{document}
